\documentclass[../main.tex]{subfiles}

\begin{document}

\chapter{Appendix B}

I wanted to include my final thoughts on these notes and notice on
some topics that were not covered here, but I guess I can use the
help of appendix.

First, as mentioned in the prefix, these are the notes that I took
while learning the essence of linear algebra in 8 weeks. That being
said, while most of the fundamental concepts were covered, other
topics like dual space, cross product, euclidean inner product,
and more advanced concepts including Markov chains and Jordan
canonical forms are not included here. After the initial release,
I will continue to update these notes for the second edition as I will
likely take linear algebra course again in college as a student who
will major in pure math. In the meantime, I highly recommend watching
3Blue1Brown series on linear algebra as it really helped me visually
understand different concepts and formulas.

As I also study different concepts in machine learning, I realized
how important linear algebra is outside ``mathematics.'' Literally,
vectors are used everywhere including, but not limited to
encoding/decoding and backpropagation. One of the reasons why I
find math so fascinating is due to duality and applications that
just seems so perfect. Matrices can be interpreted with heavy
computation while it can also be interpreted geometrically with
change in area. The two seemingly unrelated topics actually imply
the exact same concept. Similarly, what else can we use instead of
matrix multiplication and addition for forward pass in machine
learning? The way it seems so intuitive and ``perfect'' is I think
of the most beautiful side of mathematics. Really, if you completed
these notes, I highly recommend learning math behind machine learning
if like finding different applications of mathematics. More
information can be found in my
\href{null}{Notes for Machine Learning}.
\end{document}


